\section{Базовые концепции объектно-ориентированного программирования}

\textbf{Объектно-ориентированное программирование (ООП)} --- парадигма, в которой программа строится как система взаимодействующих объектов. Для объекта обычно выделяют \textbf{состояние}, \textbf{поведение} и \textbf{идентичность}. Основные понятия ООП: класс и объект, инкапсуляция, наследование, полиморфизм и абстракция.

\subsection{Объекты и классы}

\textbf{Объект} --- программная сущность, обладающая состоянием и набором операций над ним. Например, объект ``матрица'' хранит размеры и элементы и предоставляет операции сложения, умножения, вычисления нормы и т.п.

\textbf{Класс} --- описание множества объектов с общей структурой и поведением. Он задаёт поля данных, методы и правила доступа к ним; конкретный объект является \textbf{экземпляром класса}.

Следует различать класс и тип. Во многих ОО-языках класс задаёт тип, но тип может задаваться и интерфейсом без конкретного способа хранения состояния.

\subsection{Инкапсуляция}

\textbf{Инкапсуляция} --- объединение состояния объекта и операций над ним с выделением публичного интерфейса и ограничением прямого доступа к внутреннему представлению.

Например, вместо произвольного изменения внутреннего поля объекта
\begin{verbatim}
account.balance = value
\end{verbatim}
предпочтительно использовать операции
\begin{verbatim}
account.deposit(amount)
account.withdraw(amount)
\end{verbatim}
которые могут проверять допустимость изменения состояния.

Инкапсуляция позволяет поддерживать \textbf{инварианты объекта} и уменьшает зависимость внешнего кода от деталей реализации. Модификаторы \texttt{public}, \texttt{private}, \texttt{protected} являются типичными средствами сокрытия данных, но сокрытие --- лишь один из механизмов инкапсуляции.

\subsection{Наследование}

\textbf{Наследование} --- механизм построения нового класса на основе существующего. Исходный класс называют базовым, новый --- производным. Наследование естественно выражает отношение обобщения ``является'' (\textit{is-a}):
\[
\texttt{Circle} \;<:\; \texttt{Shape}.
\]
Производный класс может расширять базовый класс и \textbf{переопределять} его методы.

Различают одиночное, многоуровневое и, в некоторых языках, множественное наследование. Множественное наследование требует аккуратного разрешения неоднозначностей, например при ромбовидной иерархии.

Наследование следует отличать от отношения ``имеет'' (\textit{has-a}). Если объект является составной частью другого объекта, обычно используется \textbf{композиция}, а не наследование.

\subsection{Полиморфизм и динамическое связывание}

\textbf{Полиморфизм} --- возможность использовать общий интерфейс для объектов разных конкретных типов. Например, для набора фигур можно вызывать одну операцию \texttt{draw()} или \texttt{area()}, а выполняемая реализация определяется фактическим типом объекта.

Для классического ООП особенно важен \textbf{полиморфизм подтипов}: объект производного типа может использоваться там, где ожидается объект базового типа.

При \textbf{статическом (раннем) связывании} вызываемая функция определяется при компиляции. При \textbf{динамическом (позднем) связывании} конкретная реализация метода выбирается во время выполнения по фактическому типу объекта. В C++ этот механизм реализуется виртуальными функциями:
\begin{verbatim}
class Shape {
public:
    virtual double area() const = 0;
    virtual ~Shape() = default;
};
\end{verbatim}
Чисто виртуальная функция задаёт интерфейс без конкретной реализации; класс, содержащий хотя бы одну такую функцию, является \textbf{абстрактным}. Его объекты непосредственно не создаются, но через указатель или ссылку на базовый класс можно работать с объектами производных классов.

Следует различать \textbf{переопределение} метода в производном классе и \textbf{перегрузку} --- существование нескольких функций с одним именем и различными параметрами.

Кроме полиморфизма подтипов существуют параметрический полиморфизм (шаблоны, обобщённые типы) и ad-hoc-полиморфизм (например, перегрузка операций).

\subsection{Абстракция, интерфейсы и композиция}

\textbf{Абстракция} --- выделение существенных свойств объекта и отбрасывание деталей, несущественных для данной задачи. Пользователю численного решателя, например, достаточно знать контракт
\[
\texttt{solve}(A,b)\longrightarrow x,
\]
не зная внутренней реализации алгоритма.

\textbf{Интерфейс} задаёт набор операций, которые объект обязан поддерживать. Разные реализации одного интерфейса можно взаимозаменять, что уменьшает связанность компонентов программы.

\textbf{Композиция} означает построение сложного объекта из более простых объектов. Например,
\[
\texttt{Car has an Engine}.
\]
Она часто гибче наследования, поскольку вложенный компонент можно заменить независимо от иерархии классов. Поэтому при отсутствии естественного отношения подтипов обычно предпочитают композицию наследованию.

\subsection{Основные достоинства и ограничения}

ООП позволяет локализовать состояние, разделить программу на компоненты с ясными интерфейсами, повторно использовать реализации и расширять программу через полиморфизм. Однако чрезмерно глубокие иерархии классов и слишком большое число мелких объектов могут усложнить архитектуру и ухудшить локальность данных; в высокопроизводительных численных программах ООП поэтому часто сочетают с процедурным, обобщённым и data-oriented подходами.

\subsection{Что важно сказать на экзамене}

Нужно чётко определить \textbf{объект и класс}, затем объяснить три центральных механизма:
\[
\boxed{\text{инкапсуляция}\;\longrightarrow\;\text{наследование}\;\longrightarrow\;\text{полиморфизм}}.
\]
Для полиморфизма желательно отдельно назвать \textbf{динамическое связывание} и виртуальные функции, а также отличить наследование (\textit{is-a}) от композиции (\textit{has-a}).

\newpage
